\section{Weighted least representatives retain component identity}
\label{sec:weighted-transversal}

An orbit histogram records how many components have each additive weight.
It does not identify where these components occur in the original interval.
This distinction matters when a later geometric query must place one mark on
each boundary component, while omitting isolated interior vertices. Selecting
one point per orbit and separately computing a weight histogram loses the
association needed to make that decision.

The new \code{weighted\_orbit\_transversal} returns disjoint half-open runs
$[a,b)$ of least original representatives, each tagged by the exact weight
of its orbit. A run of length $b-a$ represents that many different orbits
with the same weight. Signed additive weights and overlapping input intervals
are supported. The run endpoints and weights are binary integers; neither
large runs nor their orbits are expanded. Adjacent runs with equal tags are
coalesced, but labels separated by a deleted interval remain separated.

\subsection{Forward transport and its invariant}

Discovery uses the monotone forward interval-orbit trace already described
in the preceding sections. Its weight transport keeps a piecewise constant
function on the current interval. Its label transport keeps the ordered
intervals of original points that are still live. At every stage, the weight
on a live point is the sum of the original weights already folded into it.
The live labels are increasing, and folding removes a point in favour of a
smaller point in the same original orbit.

Deletion, trimming, periodic merging and transmission change the description
of the equivalence relation without changing these labels or their weights.
Translation and reflection truncation add the removed points' weights to
their surviving images and retain the appropriate prefix of live labels.
A static-gap contraction finishes one orbit per point of the gap. Intersecting
the gap, weight runs and live-label intervals emits the corresponding original
labels with their final orbit weights. An ordered three-way sweep performs
this intersection without visiting individual points.

\begin{proposition}
For a complete accepted monotone trace, the emitted labels are exactly the
least original representatives, and every emitted tag is the additive weight
of the represented orbit.
\end{proposition}
\begin{proof}
The label invariant is the one used by the unweighted least-representative
operator. Weight transport preserves the sum on each equivalence class.
At a static contraction, a current point has no remaining identification with
another live point, so its weight is final. Its original label is least by
monotonicity. The complete trace removes every live point, and therefore
emits exactly one representative for each original orbit.
\end{proof}

The implementation additionally checks that the tag histogram equals the
weighted histogram, that total weight is conserved, and that the emitted
cardinality equals the ordinary orbit count. These checks supplement the
transport invariant. Histogram equality alone would accept a permutation
of tags among representatives and is therefore insufficient.

\subsection{An independent backward verifier}

The certificate contains the existing weighted interval-orbit certificate
and the proposed representative runs. Replay first authenticates the ordinary
trace and its additive weight transport. It then uses a separate construction:
record weighted emissions in the current coordinates of each contraction,
and restore the contractions backwards by inserting their removed gaps.
An interval crossing an insertion anchor is split before its labels are
lifted. Truncations insert no selected points: the removed suffix contains no
least representative. This verifier does not call the producer's forward
live-label compiler.

Only monotone forward traces are accepted. A trace using a global reversal
may preserve the histogram while changing the least-label invariant. Capped
discovery returns an inconclusive result without a selector or certificate;
external cancellation propagates. JSON transport retains large integers
through the existing exact hexadecimal encoding.

\subsection{Cost and validation scope}

Tag emission is linear in the ordered partitions traversed at that
contraction. The reference replay deliberately uses simpler backward gap
insertion, with polynomial work in the trace length, emitted-run count,
gap count and integer bit lengths. These are bounds in terms of the encoded
trace and realized weighted partitions; they do not assert a new worst-case
bound on every weighted partition generated by the scheduler.

The regression suite compares 750 seeded finite signed-pairing systems with
an independent pointwise union model and directly checks the least label and
weight of every orbit. It includes overlapping signed weights, empty systems,
a forged exchange of tags that preserves the entire histogram, cancellation,
cap exhaustion, and a $2^{16384}$-orbit example represented by one tagged run.
Producer-disabled replay accepts the large example after exact JSON transport.
